\documentclass[conference,onecolumn]{IEEEtran}
\IEEEoverridecommandlockouts
% The preceding line is only needed to identify funding in the first footnote. If that is unneeded, please comment it out.
\usepackage{cite}
\usepackage{amsmath,amssymb,amsfonts}
\usepackage{algorithmic}
\usepackage{graphicx}
\usepackage{textcomp}
\usepackage{xcolor}
\def\BibTeX{{\rm B\kern-.05em{\sc i\kern-.025em b}\kern-.08em
    T\kern-.1667em\lower.7ex\hbox{E}\kern-.125emX}}
\begin{document}

\title{Topology-Aware Mapping for Photonic Interconnects:\\
A Lightmatter Passage Case Study on TPU-Class Clusters}

\author{
\IEEEauthorblockN{Simone Surdi}
\IEEEauthorblockA{simonesurdi@gmail.com}
\and
\IEEEauthorblockN{Teammate Name}
\IEEEauthorblockA{email address}
}

\maketitle


\section{Project Overview}
This project is based on the AccelForge cluster-mapping example and extends it
with a topology-aware photonic interconnect model calibrated against the
publicly disclosed Lightmatter Passage family (EVK50, EVK100, L200/L200X, M1000).

\textbf{Problem and motivation.}
In today's accelerator clusters (e.g.\ TPU v5p's 4$\times$4$\times$4 cube and the
recently disclosed Ironwood rack), inter-chip interconnect (ICI) is increasingly
the dominant cost for distributed inference of large language models. Electrical
800G links consume on the order of 10~pJ/bit and scale poorly in 3D torus
topologies for traffic patterns such as shared-weight multicast and remote HBM
fetch. Lightmatter's Passage family claims 1.9--4.8~pJ/bit at per-fiber
bandwidths up to 57.3~Tbps, and critically supports \emph{native optical
multicast} and (for M1000) an integrated optical circuit switch.
If these links materialize, the cost model that every accelerator mapper
implicitly assumes --- nearest-neighbor traffic is cheap, remote traffic is
expensive --- may no longer hold, and mappers that do not know about the
interconnect technology will leave large amounts of energy and latency on the
table.

\textbf{Key questions.}
\begin{itemize}
    \item Under what traffic patterns (node-to-DRAM, node-to-node, shared-DRAM
    multicast) does substituting electrical ICI with photonic ICI actually
    change the \emph{energy and latency} that the AccelForge mapper reports?
    \item Does the \emph{optimal distributed-inference mapping choice} change
    when the interconnect is photonic, or do photonic links simply make the
    same mapping cheaper?
    \item How does this answer shift across Lightmatter technology generations
    (EVK50 $\rightarrow$ M1000) and across cluster topologies (2D mesh, 3D
    torus, optical crossbar, switched fabric)?
    \item What does this tell us about whether a dedicated \emph{photonic
    compiler} / topology-aware mapper would be worth building as a follow-on?
\end{itemize}

\textbf{Hypothesis.}
We expect that photonic links with native multicast will dramatically reduce
energy and latency for traffic patterns whose critical path leaves the local
HBM bank --- specifically shared-DRAM multicast and node-to-DRAM --- because
these patterns pay the electrical cost multiple times on a mesh or torus but
pay it once on a multicast-capable optical fabric. We expect node-to-node
patterns with good locality to improve only marginally, because nearest-neighbor
electrical links are already well matched to that traffic. Most importantly,
we hypothesize that for at least some workloads the optimal distributed
mapping will \emph{flip} under photonic interconnect (e.g.\ from
\texttt{node\_to\_node} to \texttt{shared\_dram\_multicast}), demonstrating
that an interconnect-agnostic mapper is strictly suboptimal once photonics are
introduced --- motivating future work on a photonic-aware compiler.


\section{Technical Contributions}
\textbf{Prior work.}
AccelForge and the Timeloop/Accelergy lineage provide analytical energy and
latency models for single-chip accelerators, but treat cluster-level
interconnect as a flat, bandwidth-limited channel. Cluster-aware modeling
exists for electrical ICI (e.g.\ TPU v4/v5p 3D-torus analyses in the course
readings) and separately there is a growing body of public material on
co-packaged optics from Lightmatter, Ayar Labs, and the OCP CPO reference
architecture, but the two have not been tied together in a mapper-visible way
for a TPU-class cluster.

\textbf{Our approach.}
We add a \emph{topology-aware network cost model} to AccelForge's iterative
mapper. Each topology (2D mesh, 3D torus, optical crossbar, photonic switched
fabric) exposes a per-traffic-pattern cost that depends on (i) the physical
technology (bandwidth and pJ/bit from a public-sourced technology catalog) and
(ii) whether the technology supports multicast. The iterative mapper loop is
extended so that network costs are fed back into AccelForge's logical-memory
action counts; the loop runs until the tile/residency signature is a fixed
point. On top of this, we add a \emph{distributed mapping selector} that, for
each (workload, architecture) pair, evaluates the candidate mappings
\texttt{node\_to\_node}, \texttt{node\_to\_dram}, and
\texttt{shared\_dram\_multicast} and reports both the best mapping and whether
it differs from the electrical baseline --- this is the core experimental
hook for the mapping-adaptation study.

\textbf{How this differs from a naive ``swap in lower pJ/bit'' study.}
Rescaling link energy alone only shows that optics reduce interconnect energy.
Our contribution is that the \emph{mapper sees} the new topology and can
therefore choose different tile shapes, residencies, and distributed mapping
strategies --- which is what a real compiler would have to do, and what
exposes qualitative (not just quantitative) changes.

\textbf{Tradeoffs and risks.}
\begin{itemize}
    \item \emph{Model fidelity.} Public Lightmatter numbers are headline
    wall-plug efficiencies; real deployments will have laser, packaging, and
    thermal overheads we abstract into a single pJ/bit figure. We mitigate by
    running sensitivity sweeps over pJ/bit and bandwidth.
    \item \emph{Multicast semantics.} Native optical multicast is modeled as a
    one-hop broadcast; in practice it depends on the switch configuration and
    fan-out limits. The model is deliberately optimistic and is labeled as
    such in the reported plots.
    \item \emph{AccelForge fixed-point convergence.} The iterative mapper may
    oscillate or converge slowly when interconnect costs change drastically;
    we report convergence traces as a first-class result.
\end{itemize}

\textbf{Why we think it will work.}
A preliminary implementation on a 16-node cluster already shows that
\texttt{shared\_dram\_multicast} traffic sees on the order of 70\% energy
reduction with EVK100-class links, and that for \texttt{prefill\_projection} on
a 4$\times$4$\times$4 TPU cube the optimal distributed mapping changes from
\texttt{node\_to\_node} (electrical) to \texttt{shared\_dram\_multicast}
(photonic) --- which is direct evidence that the hypothesis holds in at least
one concrete case.

\textbf{Key challenges.}
Keeping the technology catalog faithful to public sources without overfitting
to marketing numbers; keeping the iterative mapper numerically stable when
interconnect costs drop by an order of magnitude; and cleanly separating
``bandwidth wins'' from ``multicast wins'' from ``mapping wins'' in the
reported metrics.


\section{Evaluation}
\textbf{Baseline and fair comparison.}
The baseline is an electrical cluster using 800~Gbps per-direction ICI at
10~pJ/bit --- chosen because it matches the TPU v5p ICI operating point
disclosed in public Google Cloud documentation. We hold the compute fabric
constant (identical TPU-like MAC array, HBM, and Unified Buffer defined in
\texttt{arch/tpu\_lightmatter\_cluster.yaml}) across all experiments; only the
interconnect layer and the mapper's knowledge of it are varied. Where area is
relevant (switched-fabric variants), we report per-chip interconnect
bandwidth density (Tbps/mm$^2$) alongside energy and latency so that
higher-bandwidth links are not implicitly credited for spending more area.

\textbf{Experiment 1 --- 16-node cluster study: topology swap.}
\emph{Class concept:} interconnect energy vs.\ compute/memory energy; traffic
pattern vs.\ topology.
\emph{Setup:} Three distributed-inference traffic patterns
(\texttt{node\_to\_dram}, \texttt{node\_to\_node},
\texttt{shared\_dram\_multicast}) $\times$ three representative LLM workloads
(\texttt{decode\_mlp}, \texttt{kv\_attention}, \texttt{prefill\_projection})
$\times$ two topologies (electrical 4$\times$4 mesh at 10~pJ/bit vs.\ Lightmatter
EVK100 optical crossbar). Metrics: per-workload energy (J) and latency (s).
\emph{Communication:} Grouped bar chart, one group per workload, bars for each
(topology, scenario) combination, with a companion CSV
(\texttt{cluster\_improvements.csv}) reporting percent reductions.
\emph{Expected trend:} Largest reductions on multicast and DRAM-fetch patterns
(70\%+ energy), near-zero on pure nearest-neighbor node-to-node --- validating
that photonics help specifically where the electrical mesh is forced to pay
hop-count or fan-out overhead.

\textbf{Experiment 2 --- TPU cube ICI replacement: technology generations.}
\emph{Class concept:} scaling interconnect bandwidth and pJ/bit; technology
generations.
\emph{Setup:} A 64-chip 4$\times$4$\times$4 TPU cube (matching v5p / Ironwood
building-block geometry) with the electrical torus baseline replaced by five
photonic variants from the technology catalog (EVK100 torus, L200 switched,
L200X switched, M1000 switched, plus the EVK50 reference point). Three traffic
archetypes are run: \texttt{all\_to\_all}, \texttt{weight\_refill},
\texttt{multicast}. Metrics: cube-level energy (J), tail latency (s), and
percent reduction vs.\ the electrical torus.
\emph{Communication:} Heatmap (workload $\times$ variant) of percent energy
reduction; line plot of reduction vs.\ technology generation to visualize the
bandwidth/efficiency Pareto.
\emph{Expected trend:} Monotonically better energy as we move from EVK50 to
M1000, with the largest wins on multicast patterns (approaching two orders of
magnitude reduction because the electrical torus pays the fan-out in
nearest-neighbor hops). Latency gains should saturate earlier than energy
gains because the torus is already latency-competitive for short messages.

\textbf{Experiment 3 --- Mapping-choice adaptation.}
\emph{Class concept:} the mapper as a joint tile/residency/placement
optimizer; how cost-model changes propagate into mapping decisions.
\emph{Setup:} For each (workload, architecture) pair in Experiments 1 and 2,
run the AccelForge iterative mapper to a fixed point under each candidate
distributed mapping (\texttt{node\_to\_node}, \texttt{node\_to\_dram},
\texttt{shared\_dram\_multicast}), and record (a) the fixed-point tile /
residency signature and (b) which distributed mapping is optimal. Compare
against the electrical baseline to identify cases where the optimal mapping
\emph{changes} under photonics.
\emph{Metrics:} Fraction of (workload, architecture) cases where the optimal
distributed mapping differs from the electrical baseline; mapping-aware
energy gain on the 4$\times$4$\times$4 cube (in addition to the raw
interconnect-swap gain).
\emph{Communication:} Table (\texttt{optimal\_mapping\_choices.csv}) with one
row per (workload, architecture) listing the optimal mapping and a
``Changed vs.\ baseline'' flag, plus a bar chart decomposing the total
improvement into ``interconnect swap'' vs.\ ``mapping change'' contributions.
\emph{Expected trend:} At least one workload should change its optimal
distributed mapping under photonics, contributing additional percent-level
energy improvement beyond the raw link replacement --- which would be the
direct motivation for a future photonic-aware compiler. We expect
\texttt{prefill\_projection} in particular to flip from \texttt{node\_to\_node}
to \texttt{shared\_dram\_multicast} because its weight matrices are large
enough that multicast amortizes over many consumers once multicast is cheap.

\textbf{Extended sensitivity analysis.}
Beyond the three primary experiments, we sweep (i) per-link bandwidth from
0.8--57.3~Tbps and (ii) pJ/bit from 1.9--10.0 to decouple ``bandwidth wins''
from ``efficiency wins,'' and we run a node-count scaling sweep from 2 to 128
nodes to show how the advantage evolves with cluster size.


\section{Timeline and Work Distribution}
\textbf{Weeks 1--2 --- Network model (M1).}
Build the topology-aware cost model (\texttt{Topology}, \texttt{Mesh2D},
\texttt{Torus3D}, \texttt{OpticalCrossbarTopology}) and wire it into
AccelForge's logical-memory action counts. \emph{Owner:} Teammate A,
reviewed by Teammate B.

\textbf{Weeks 2--3 --- Iterative mapper loop (M2).}
Implement \texttt{solve\_iteratively} and \texttt{compare\_topologies};
instrument fixed-point convergence tracing. \emph{Owner:} Teammate B,
reviewed by Teammate A.

\textbf{Weeks 3--4 --- Proposed optical topology (M3) + technology catalog.}
Implement the 16-node optical crossbar study and build the public-sources
technology catalog (EVK50, EVK100, L200, L200X, M1000). \emph{Both teammates,
split by topology type.}

\textbf{Weeks 4--5 --- TPU cube ICI study and mapping adaptation.}
Run the 4$\times$4$\times$4 cube study, implement the distributed-mapping
selector, and generate the mapping-change table. \emph{Teammate A leads the
cube architecture; Teammate B leads the mapping-adaptation analysis.}

\textbf{Week 5 --- Sensitivity, presentation assets, and writeup.}
Bandwidth and pJ/bit sweeps, node-count scaling, generation of final plots
and CSVs under \texttt{results/presentation/}, plus this report. \emph{Both
teammates.}

\textbf{Stretch (post-course) --- Photonic compiler.}
If time permits, prototype a mapper pass that takes the interconnect
technology descriptor as a first-class input and emits a mapping chosen from
a photonic-aware search space. This is flagged as future work and motivates
the broader research agenda.

\end{document}
